\subsection{10 марта. Спектр. Резольвента}

Ситуация: $E(\C)$~--- БП, $A \in \mathcal{L}(E)$.

\begin{definition}
Назовем $\lambda$~--- регулярное значение, или $\lambda \in \rho(A)$, если уравнение
\begin{equation}
    \label{equation_regular_val}
    (A - \lambda I) x = y
\end{equation}
имеет единственное решение $\forall y \in E$, которое непрерывно зависит от $y$. То есть

$$
\exists A_\lambda^{-1} \in \mathcal{L}(E).
$$

\end{definition}

\begin{definition}
    Спектр: $\sigma(A) = \C \setminus \rho(A)$.
\end{definition}

\begin{definition}
    Резольвента: $\lambda \in \rho(A), \; R_\lambda = A_\lambda^{-1}$.
\end{definition}

\begin{statement}[$-1$]
    Если $|\lambda|>\|A\|$, то $\lambda \in \rho(A)$. \\
    Смысл: спектр сидит где-то внутри круга радиуса нормы оператора.
\end{statement}

\begin{proof}
    \begin{align*}
        (A-\lambda I)x &= y,\;\; \lambda \neq 0,\\
        (I - \frac{1}{\lambda}A)x &= -\frac{1}{\lambda} y.
    \end{align*}

    По теореме~\nameref{theorem8.2}. $\|-\frac{1}{\lambda}A\| < 1$.
    \[(I-\frac{1}{\lambda}A)^{-1}=\sum_{n=0}^\infty \frac{1}{\lambda^n}A^n\]
\end{proof}

\begin{theorem}[10.1]
\begin{enumerate}
    \item Спектр --- непустое замкнутое множество
    \item Спектральный ряд $r(A) = \sup_{\lambda \in \sigma} |\lambda| = r_{\text{сх}} \text{ ряда Неймана \hypertarget{thm10.1-neihman}{(**)}} = \lim\limits_{n \to \infty} \sqrt[n]{\|A^n\|}$.
\end{enumerate}
\end{theorem}

\textbf{Идея:} рассматриваются функции $F: \C \to \L(E)$. Неудобство: нет коммутативности умножения для операторов

\textbf{Немного ТФКП для пространства операторов:} непрерывность и дифференцируемость вводится по определению. Вводим интеграл (соответственно появляется теорема Коши), ряд Тейлора и Лорана, теорема Лиувилля.... (подробнее см. Колмогоров-Фомин, приложение <<Банахова алгебра>>)

\begin{statement}[$0$]
    $\rho(A)$ - открытое множество
\end{statement}

\begin{proof}
    \textbf{Идея:} теорема~\nameref{theorem8.3}.
    $\lambda_0 \in \rho(A) \; \exists R_{\lambda_0}=A_{\lambda_0}^{-1}$
    $$\frac{()}{()} \in L(E)$$.
    % TODO ДОПИСАТЬ
\end{proof}

\begin{corollary}
$\sigma(A)$~--- замкнутое.
\end{corollary}

Возьмем какое-то $\lambda \in \C$. Есть два вариант: либо ядро $A_\lambda$ тривиально, либо нетривиально (тогда это собственное значение, $\lambda \in \sigma_p(A)$~--- точечный спектр). Случай тривиального ядра разделяется на 3
\begin{enumerate}
    \item $\im A_\lambda=E$. Тогда по теоерме~\nameref{theorem8.4} $\exists A_\lambda^{-1}$ ($\rho(A)$)
    \item $\im A_\lambda \neq E$, $\overline{\im A_\lambda} = E$ ($\sigma_C(A)$)
    \item $\overline{\im A_\lambda} \neq E$ ($\sigma_R(A)$~--- остаточный спектр)
\end{enumerate}

\textbf{Итого:} вся комплексная плоскость делится на $\rho(A), \sigma_p(A), \sigma_C(A), \sigma_R(A)$.

Рассмотрим произведение:
\begin{align*}
    (Ax, Ax) &= (x, A^*Ax) = (\sum\limits_{i = 1}^n \xi_i e_i, \sum\limits_{i = 1}^n \lambda_i \xi_i e_i )\\
    &= \sum\limits_{i = 1}^n \lambda_i \xi_i^2 \leq (\max\limits_{i} \lambda_i) \sum\limits_{i = 1}^n \xi_i^2 = (\max\limits_{i} \lambda_i) \|x\|^2.
\end{align*}
где $x = \sum\limits_{i=1}^n \xi_i e_i$. Получили оценку на норму $A$.

\begin{theorem}[Гильберта-Шмидта]\label{hilb-shmidt}
    Если $H$ - сеперабельное гильбертово пространство, то в $H$ существует ОНБ, ссостоящий из собственных векторов оператора $A$, где $A$ - компактный самосопряженный оператор.
\end{theorem}

\begin{statement}[$1$]
    $R_\lambda$ - непрерывная функция на $\rho(A)$
\end{statement}

\begin{proof}
    Если взять некую точку $\lambda_0 \in \rho(A)$, то она входит в $\rho(A)$ с некоторой окрестностью по прошлому утверждению.

    В теореме~\nameref{theorem8.3}:
    $$\|(A+\Delta A)^{-1}-A^{-1}\| \leqslant \frac{\|A^{-1}\|^2 \|\Delta A \|}{1-\|A^{-1}\| \|\Delta A\|}$$
    Подставляем $R_\lambda$:
    $$\frac{\|R_{\lambda_0} \|^2 |\Delta \lambda\|}{1 - \| R_{\lambda_0}\| \|\Delta \lambda|} \to 0$$
\end{proof}

\begin{statement}[$2$]
   Равенство Гильберта: $\lambda_0, \lambda \in \rho(A)$:
   $$R_\lambda - R_{\lambda_0} = (\lambda - \lambda_0)R_\lambda R_{\lambda_0}.$$
\end{statement}

\begin{proof}
\[
    R_\lambda - R_{\lambda_0} = \\
    R_\lambda(A_{\lambda_0} R_{\lambda_0})-(R_\lambda A_\lambda)R_{\lambda_0}=R_\lambda(A_{\lambda_0}-A_\lambda)R_{\lambda_0}.
\]
\end{proof}

\begin{statement}[$3$]
   Функция $R_\lambda$ дифф. на области $\rho(A)$.
\end{statement}

\begin{proof}
    По равенству Гильберта получаем
    \begin{align*}
        \frac{R_\lambda - R_{\lambda_0}}{\lambda - \lambda_0} =_{\text{утв 2}} \frac{(\lambda - \lambda_0)R_\lambda R_{\lambda_0}}{(\lambda - \lambda_0)} \to_{\text{утв 1}} R_{\lambda_0}^2.
    \end{align*}
\end{proof}

\begin{statement}[$4$]
    Радиус сходимости ряда~\hyperlink{thm10.1-neihman}{**} равен $r(A)$ и равен $\overline{\lim_{n \to \infty}} \sqrt[n]{\|A^n\|}$
\end{statement}

\begin{proof}
\begin{enumerate}
    \item $|\lambda_0| > r(A) \implies \lambda_0 \in \rho(A)$. $R_\lambda$ раскладывается в ряд Лорана единственным образом, значит это ряд Неймана $\implies r_{\text{сх}} \leqslant r(A)$
    \item $|\lambda_0| < r(A)$. Утверждаем, что ряд Неймана~\hyperlink{thm10.1-neihman}{**} расходится в точке $\lambda_0$. Если бы он сходился в этой точке, то он бы сходился и при любом $\lambda: |\lambda| > |\lambda_0|$. Получаем, что $r_{\text{сх}} \geqslant r(A)$.
\end{enumerate}
    
\end{proof}

\begin{statement}[$5$]
    Спектр непустой
\end{statement}

\begin{proof}
    Предположим противное. то есть $\forall \lambda \in \C \implies \lambda \in \rho(A) = \C$.

    \hyperlink{thm10.1-neihman}{**} $\implies \|R_\lambda\| \leq \frac{1}{|\lambda|} \frac{1}{1-\|A\|/|\lambda|}=\frac{1}{|\lambda| - \|A\|} \to 0, \lambda \to \infty$

    Пришли к противоречию
\end{proof}

\begin{statement}[$6$]
    $\lambda \in \sigma(A) \implies \lambda^n \in \sigma(A^n)$
\end{statement}

\begin{proof}
    Предположим противное: $\lambda^n \not\in \sigma(A^n)$, то есть $\lambda^n \in \rho(A^n) \sim (A^n-\lambda^n I)^{-1} \in \L(E)$
    \[(A^n-\lambda^nI)=(A-\lambda I)(A^{n-1} + \ldots + \lambda^{n-1}I)\]

    Умножим справа на $A^n-\lambda^n I)^{-1}$ получим правый обратный, слева - левый обратный
\end{proof}

\begin{statement}[$7$]
    $r(A)=\lim \sqrt[n]{\|A^n\|}$
\end{statement}

\begin{proof}
    Рассмотрим соотношение $$(r(A))^n \leq \left[r(A^n)\right]^{\frac1n} \leq (\|A^n\|)^{\frac{1}{n}}$$
    которое верно в силу прошлого утверждения.

    По пункту 4: $r(A)=\overline{\lim\limits_{n \to \infty}} \sqrt[n]{\|A^n\|}$, и видно, что эти числа совпадают.

    $$\alpha_n=\sqrt[n]{\|A^n\|}$$
    $$\overline{\lim} \alpha_n \leq \alpha_n \; \forall n$$
    $$\overline{\lim} \alpha_n \leqslant \underline{\lim} \alpha_n$$
\end{proof}

\begin{exercise}
    $l_2(\C), \{\lambda_n\}$ - ограниченная, $Ax=(\lambda_1 x_1, \lambda_2 x_2, \ldots)$. Док-ть $\sigma(A)=\overline{\{ \lambda_n \}}$
\end{exercise}